% Part: turing-machines 
% Chapter: machines-computations 
% Section: introduction

\documentclass[../../../include/open-logic-section]{subfiles}

\begin{document}

\olfileid{tur}{mac}{int}
\olsection{Pambuka}

Apa tegese fungsi, upamane saka $\Nat$ menyang $\Nat$,
iku \emph{komputabel}? Salah siji wangsulan kang paling
awal lan paling kawentar yaiku fungsi komputabel yen bisa
diitung dening mesin Turing. Pangerten iki diandharake
dening Alan Turing ing taun 1936. Mesin Turing minangka
tuladha \emph{model komputasi}---cara kang cetha sacara
matematis kanggo netepake gagasan ``prosedur komputasi''.
Tegese gagasan iki kanthi persis isih dadi rembugan,
nanging akeh kang sarujuk yen mesin Turing minangka
salah siji cara kanggo nemtokake prosedur komputasi.
Sanajan istilah ``mesin Turing'' bisa nuwuhake gambaran
mesin fisik kanthi bagean-bagean kang obah, sejatine
mesin Turing mung konstruksi matematis, mula dadi
idealisasi prosedur komputasi. Tuladhane, kita ora
matesi wektu utawa memori kang dibutuhake mesin Turing:
mesin Turing bisa ngitung sawijining perkara sanajan
komputasine mbutuhake papan panyimpenan utawa langkah
kang luwih akeh tinimbang cacah atom ing alam semesta.

\begin{explain}
Mbokmenawa luwih becik mesin Turing dianggep
minangka program kanggo sawijining mekanisme
khayalan kang mligi. Mekanisme iki dumadi saka
\emph{pita} lan \emph{sirah maca-nulis}. Ing versi
mesin Turing kang kita gunakake, pita iku tanpa
wates ing siji arah (menyang tengen), lan kabagi
dadi \emph{kothak-kothak}; saben kothak bisa ngemot
simbol saka \emph{alfabet} winates. Alfabet kaya
mangkono bisa ngemot simbol beda-beda kanthi cacah
sembarang, nanging kita biasane cukup nganggo telu:
$\TMendtape$, $\TMblank$, lan $\TMstroke$. Nalika
mekanisme diwiwiti, pita kosong (yaiku, saben
kothak ngemot simbol $\TMblank$) kajaba kothak
paling kiwa, kang ngemot $\TMendtape$, lan
sawetara kothak winates kang ngemot \emph{input}.
Ing saben wektu, mekanisme ana ing salah siji
saka \emph{kahanan} kang cacahe winates. Ing wiwitan,
sirah maca kothak paling kiwa lan mekanisme ana
ing \emph{kahanan wiwitan} kang wis ditemtokake.
Ing saben langkah lumakune mekanisme, isi kothak
kang lagi diwaca, kahanan mekanisme, lan program
mesin Turing nemtokake apa kang kedadeyan sabanjure.
Program mesin Turing diwenehake dening fungsi
parsial kang nampa kahanan~$q$ lan simbol~$\sigma$
minangka input lan ngasilake tripel~$\tuple{q', \sigma',
  D}$. Nalika mekanisme ana ing kahanan $q$ lan
maca simbol $\sigma$, mekanisme ngganti simbol
ing kothak saiki dadi $\sigma'$, sirah obah
menyang kiwa, menyang tengen, utawa tetep ing
panggonane manut apa $D$ iku $\TMleft$,
$\TMright$, utawa $\TMstay$, banjur mekanisme
mlebu kahanan~$q'$.

Tuladhane, gatekna kahanan ing \olref{fig:tm}.
\begin{figure}
  \olasset{\olpath/assets/diagrams/turing-machine.tikz}
  \caption{Mesin Turing kang nglakokake programe.}
  \ollabel{fig:tm}
\end{figure}
Bagean pita mesin Turing kang katon ngemot
simbol pucuk pita $\TMendtape$ ing kothak paling
kiwa, banjur telung simbol $1$, siji $0$,
lan patang simbol $1$ maneh. Sirah lagi maca
kothak katelu saka kiwa, kang ngemot
simbol~$1$, lan ana ing kahanan~$q_1$---kita
ngarani ``mesin lagi maca $1$ ing
kahanan~$q_1$.'' Yen program mesin Turing
kanggo input $\tuple{q_1, 1}$ mbalekake
tripel $\tuple{q_2,
0, \TMstay}$, mesin banjur ngganti~$1$ ing
kothak katelu dadi~$0$, ora ngowahi panggonan
sirah maca-nulis, lan ngalih menyang kahanan~$q_2$.
Yen banjur program mbalekake $\tuple{q_3, 0, \TMright}$
kanggo input $\tuple{q_2, 0}$, mesin banjur
nimpa~$0$ nganggo~$0$ maneh (mula isi pita ing
sangisore sirah maca-nulis ora owah), obah
siji kothak menyang tengen, lan mlebu
kahanan~$q_3$. Mangkono sateruse.

Kita ngarani mesin \emph{mandheg} nalika
nemoni kahanan, $q_n$, lan simbol, $\sigma$
kang ora nduwé instruksi kanggo
$\tuple{q_n, \sigma}$, yaiku fungsi transisi
ora ditegesi kanggo input $\tuple{q_n,\sigma}$.
Kanthi tembung liya, mesin ora nduwé instruksi
kang bisa ditindakake lan banjur mandheg
makarya. Kadhang kala mandheg iku diwakili
dening kahanan mandheg mligi~$h$. Iki bakal
dituduhake luwih rinci ing bagean sabanjure.
\end{explain}

\begin{digress}
Kaistimewan tulisan Turing, ``On computable numbers,''
yaiku ora mung menehi definisi formal, nanging
uga argumen yen definisi iku nggambarake
pangerten komputabilitas miturut intuisi.
Saka definisine, kudune cetha yen fungsi kang
bisa diitung mesin Turing uga komputabel miturut
pangerten intuisi. Turing menehi telung jinis
argumen kanggo arah kosok baline, yaiku yen
saben fungsi kang lumrahe kita anggep komputabel
uga bisa diitung dening mesin kaya mangkono.
Telung argumen iku (miturut tembunge Turing):
\begin{enumerate}
\item Ngrujuk langsung marang intuisi.
\item Bukti yen rong definisi iku ekuivalen
  (yen definisi anyar luwih gampang ditampa
  miturut intuisi).
\item Menehi tuladha kelas-kelas gedhe
  saka wilangan kang komputabel.
\end{enumerate}
Tujuan kita yaiku nyoba netepake pangerten
komputabilitas ``miturut prinsip'', yaiku
tanpa nggatekake watesan wektu lan papan
kang nyata. Mesthi wae, sawise definisi
komputabilitas kang paling amba ditemtokake,
kita bisa nerusake ngrembug komputasi kanthi
sumber daya winates; iki dadi inti saka
bidang ``kompleksitas komputasional''.
\end{digress}

\begin{history}
Alan Turing ngripta mesin Turing ing taun 1936.
Sanajan nalika semana kang dadi kawigatene
yaiku bisa-orane logika orde kapisan diputusake,
tulisane diarani tulisan kang wigati banget
bab dhasar rancangan komputer. Ing tulisan
iku, Turing nggatekake wilangan real komputabel,
yaiku wilangan real kang pangembangane ing
sistem desimal bisa diitung; nanging dheweke
nyebut yen pangerten iku gampang dicocogake
kanggo fungsi komputabel ing wilangan asli,
lan sateruse. Gatekna yen iki limang taun
sadurunge komputer serbaguna kang bisa makarya
kapisan dibangun ing taun 1941 (dening
Konrad Zuse saka Jerman ing ruang tamu
wong tuwane), pitung taun sadurunge Turing
lan kanca-kancane ing Bletchley Park mbangun
Colossus kanggo mecah kode (1943), sangang
taun sadurunge ENIAC saka Amerika (1945),
rolas taun sadurunge komputer serbaguna
kapisan saka Britania---Manchester Small-Scale
Experimental Machine---dibangun ing Manchester
(1948), lan telulas taun sadurunge wong
Amerika nyoba BINAC kanggo kapisan (1949).
Manchester SSEM minangka komputer kapisan
kang nyimpen program---mesin sadurunge
kudu disambung maneh nganggo tangan
kanggo saben tugas anyar.
\end{history}

\end{document}
